
\include{./def/yf-formatting}
\include{./def/yf-def}

\usepackage{wasysym}
%\usepackage{times}

\title{}
\author{}
\date{}


%\def\A{\mathcal{A}}
%\def\B{\mathcal{B}}
\def\C{\mathcal{C}}
%\def\D{\mathcal{D}}
\def\E{\mathcal{E}}
\def\FF{\EuScript{F}}
\def\G{\mathcal{G}}
\def\GG{\EuScript{G}}
\def\H{\mathcal{H}}
\def\I{\mathcal{I}}
\def\J{\mathcal{J}}
\def\L{\mathcal{L}}
\def\Q{\mathcal{Q}}
\def\S{\mathcal{S}}
\def\T{\mathcal{T}}
\def\V{\mathcal{V}}
%\def\W{\EuScript{W}}
\def\W{\mathscr{W}}
\def\X{\mathcal{X}}
\def\Y{\mathcal{Y}}
\def\Z{\mathcal{Z}}
\def\att{\mit{diffatt}}
\def\agm{\mathit{AGM}}
\def\boldatt{\mathbf{att}}
\def\ext{\mathit{ext}}
\def\In{\mathrm{IN}}
\def\join{\mathit{join}}
\def\Out{\mathrm{OUT}}
\def\schema{\mathit{schema}}

\def\extraspacing{\vspace{4mm} \noindent}
\def\vgap{\vspace{4mm}}
\def\vslit{\vspace{2mm}}

\begin{document}

\boxminipg{\linewidth}{
    \begin{center}
        \vspace{3mm}
    {\Large
    (CE7456 Lecture Notes 13) The External Memory Model} \\[2mm]
    {\large Yufei Tao}
    \vspace{3mm}
    \end{center}
}


So far, all our algorithm analysis has been conducted in the RAM (random access machine) model. This model is suitable for studying algorithms that aim to minimize CPU time. However, in many practical applications --- such as database systems --- data sets are too large to fit in main memory and must instead be stored on disk. Algorithms operating on such data perform numerous disk I/O operations, which are orders of magnitude slower than CPU computations and often become the main performance bottleneck. The RAM model is no longer an appropriate computation model for studying algorithms that are designed to minimize I/O cost.

\vgap

This lecture introduces the {\em external memory} (EM) model, which is the de facto model for analyzing I/O-efficient algorithms. We will then illustrate the characteristics of the model through several representative algorithms.

\extraspacing {\bf Math Conventions.} Given an integer $x \ge 1$, we use $[x]$ to represent the set $\set{1, 2, ..., x}$.

\section{The External Memory Model} \label{sec:em}

In the EM model, a machine is equipped with a memory of $M$ words and a disk of unbounded size. The disk is formatted into disjoint {\em blocks}, each having $B$ words, where $B$ is called the {\em block size}. The values $M$ and $B$ are parameters of the model satisfying $M \ge 3B$ (they cannot be regarded as constants).

\vgap

Every disk block has a unique integer address. Given such an address, a {\em read I/O operation} loads the contents of the corresponding disk block into memory, replacing the $B$ (memory) words previously stored there. Conversely, given a disk address, a {\em write I/O operation} writes $B$ memory words into the corresponding block, overwriting the $B$ words originally in that block. Henceforth, we use ``I/O'' to refer to either a read I/O or a write I/O.

\vgap

The {\em cost} of an algorithm is measured as the number of I/Os performed. CPU computation is for free but can be performed only on memory-resident data. The {\em space} of a data structure is measured as the number of disk blocks occupied.

\section{Sorting} \label{sec:sort}

Let $S$ be a set of $n \ge B$ integers, each of which fits in a word. Initially, $S$ is stored in an array on disk, which occupies $\lc n/B \rc$ blocks (that is, $B$ integers per block except the last block). The goal of the {\em sorting problem} is to produce an array of $n$ words on disk, where the elements of $S$ are arranged in ascending order.

\vgap

We show that the problem can be solved in $O(\fr{n}{B} \log_{M/B} \fr{n}{B})$ I/Os, which is known to be worst-case optimal \cite{av88}.

\subsection{Merging} \label{sec:sort:merge}

Define
\myeqn{
    m = M/B. \nn
}
The value $m$ can be regarded as the number of blocks (each with $B$ words) that the memory can accommodate.

\vgap

We first tackle the {\em merging problem} defined as follows. Let $S_1, S_2, ..., S_f$ be $f$ disjoint sets of integers, where $f \le m - 1$. Define $n = \sum_{i=1}^f |S_i|$. For each $i \in [f]$, the set $S_i$ has been sorted in an array of $|S_i|$ words stored on disk. The goal is to produce an array of $n$ words on disk, where the elements of $S_1 \cup S_2 \cup ... \cup S_f$ are arranged in ascending order.

\vgap

For each $i \in [f]$, we allocate a memory block (i.e., $B$ memory words) to serve as the {\em input buffer} for $S_i$. The $f$ input buffers use $f B \le (m - 1) B = M - B$ words of memory, leaving at least $B$ memory still available. We allocate $B$ of those words to serve as the {\em output buffer}.

\vgap

In the beginning, all buffers are empty and all sets $S_1, ..., S_f$ are marked as {\em active}. For each $i \in [f]$, the algorithm starts by loading the first block of (the array) $S_i$ into its input buffer. As long as all active sets' input buffers are non-empty, we move the smallest of the elements currently in the input buffers to the output buffer (e.g., if the element $e$ is from $S_1$, then $e$ is removed from the input buffer of $S_1$ and added to the output buffer). Whenever the output buffer is full, we perform an I/O to write the (sorted) elements there to a disk block at the end of the output array. Whenever an input buffer of an active $S_i$ (for some $i \in [f]$) is empty, we perform an I/O to read the next block of $S_i$ into the input buffer. If no such block exists (i.e., the entire $S_i$ has been read), we mark $S_i$ as {\em inactive}.

\vgap

The algorithm reads every block of each $S_i$ ($i \in [f]$) once, which requires $\lc |S_i|/B \rc$ I/Os; note that the ceiling here is necessary because $|S_i|$ may be smaller than $B$. Thus, reading all of $S_1, ..., S_f$ costs $\sum_{i=1}^f \lc |S_i|/B \rc < \sum_{i=1}^f 1 + |S_i|/B = f + n/B$ I/Os. Each block in the output (sorted) array is written once. Hence, the algorithm performs $\lc n/B \rc$ write I/Os in total. The total I/O cost is therefore $O(f + n/B)$.

\subsection{Reducing Sorting to Merging} \label{sec:sort:sort}

Next, we show how to sort a set $S$ of $n$ integers, given in a disk array of $n$ words, by using a merging algorithm as a black box.
Henceforth, we use the term {\em run} to refer to a sorted array containing a subset of the elements in $S$. Our goal is to produce one single run containing the entire $S$.

\extraspacing {\bf Initial Pass.}



\section{Remark}

\bibliographystyle{plainurl}% the mandatory bibstyle
\bibliography{ref}



\end{document}
